% EDITION19_COLLATION: EGA5_francais.pdf physical pp.60--63;
% EGAVaiello.pdf physical/printed pp.41--43. Entire section 11.
% Replaces section 11 only in v18_first_8_14_onward.tex. Section 12 is unchanged.
\Needspace{10\baselineskip}
\egasection{§11\quad Grassmann 概形}
\source{法文手稿物理第60--63页；对读\texttt{EGAVaiello.pdf}第41--43页。第十九版逐段回校}

现在我们经常要使用 $P$ 中的线性子簇，而不再只限于相对维数为0的情形或超平面的情形。\cmt{D19-11-01}{英译第41页把原来的两种相对维数写为 $0$ 和 $n-1$；法文第60页相应字母的打字字形不够清楚，形似 $r$。这里按上下文写出点与超平面这两种几何情形，不把该字母与下面重新引入的商模秩 $n$、以及后面的 $E$ 的秩 $r$ 强行等同。法文标题后另有划去的副题字样，未可靠辨定，不恢复为有效标题。}因此，显然需要引入关于 Grassmann 概形的一些记号，以及关于线性子簇构造的若干属于“初等几何”的连锁事实〔原词：\textit{sorites}〕\ednote{这个词最好怎样翻译？〔英文底本脚注57。〕}；它们全都应当放在本段开头。此外，实际中有时也要取任意线性截面，而不只取超平面截面，所以宜以这种较广的观点重新审视前面所有编号小节。

设 $E$ 是预概形 $S$ 上的拟凝聚模，$n>0$ 为整数。考虑函子
\[
(\mathbf{Sch}/S)^{\mathrm{op}}\longrightarrow\mathbf{Ens},
\]
其在 $S'$ 上的值 $\operatorname{Grass}_n(E)(S')$，为 $E_{S'}$ 的局部自由、秩为 $n$ 的商模所成的集合。\transnote{这里的商模按相容的商映射同构类计；下标 $S'$ 表示沿 $S'\to S$ 的拉回。法文写 $(\mathrm{Sch})^0/S$ 及 $E_{S'}$；英译拉回下标的排印有误。正文将反变范畴和拉回记号写清。}

这个函子可表；代表它的 $S$-预概形也记作 $\operatorname{Grass}_n(E)$。为证明可表性，考虑函子间的自然态射
\[
\operatorname{Grass}_n(E)\longrightarrow
\operatorname{Grass}_1(\textstyle\bigwedge^nE)
=\mathbf P(\textstyle\bigwedge^nE),
\]
它把 $E_{S'}$ 的每个局部自由、秩为 $n$ 的商模 $G$，对应到局部自由、秩为一的模 $\bigwedge^nG$，将后者看作 $(\bigwedge^nE)_{S'}$ 的商模。\cmt{D19-11-02}{法文第60页已经明确写出 $\bigwedge^nG$ 及源 $\bigwedge^nE_{S'}$；英文第41页在这一解释句中漏掉外幂，误把秩 $n$ 的 $G$ 仍称作秩一模。正文依照法文恢复。旧注D17-E04针对的是英译这一遗漏，不能据此声称法文手稿同样漏写。}
与 Cartan 讨论班\ednote{请把参考文献写得更精确。〔英文底本脚注58。〕}中的证明一样，可以证明这个态射“可由闭浸入表示”；于是 $\operatorname{Grass}_n(E)$ 是 $\mathbf P(\bigwedge^nE)$ 的闭子预概形。特别地，它在 $S$ 上分离；当 $E$ 为有限型时，它在 $S$ 上拟紧并且射影。\cmt{D19-11-03}{法文第61页与英文第41页都把“在 $S$ 上拟紧”无条件置于“若 $E$ 为有限型则射影”之前。一般拟凝聚 $E$ 不保证拟紧性：取 $S=\Spec k$、$E=\bigoplus_{i\geq0}ke_i$、$n=1$，则 $\mathbf P(E)=\operatorname{Proj}k[T_0,T_1,\ldots]$ 的标准开覆盖 $D_+(T_i)$ 没有有限子覆盖。正文保留旧注D17-E05采用的有限型限制；这是原断言的实际订正。}
若 $E$ 为有限表示，则 $\operatorname{Grass}_n(E)$ 在 $S$ 上也为有限表示。事实上，可以假设 $S$ 仿射，写作 $S=\Spec A$；于是 $E$ 来自 $A$ 的某个有限型子环上的一个有限型模，而 $\operatorname{Grass}_n(E)$ 的构造显然与 $S$ 上的基变换相容。\cmt{D19-11-04}{法文与英译这里都说“有限型子环上的有限型模”。这项归约可以成立：取定义 $E$ 的有限表示矩阵中的有限个系数所生成的 $\mathbf Z$-子代数 $A_0\subset A$，并用同一矩阵定义 $A_0$-模。$A_0$ 是诺特环，其有限型模本来就有限表示。因此这里不应把“有限型模”判作原稿必须改正的错误；旧中文对有限矩阵下降的解释保留于核查报告，并撤销其“原文信息不足”的归责。这里也不是断言任意有限型 $A$-模的全部关系均可如此下降。}

若 $E$ 局部自由且秩有限，则 $\operatorname{Grass}_n(E)$ 在 $S$ 上光滑；在秩至少为 $n$ 的部分，其纤维几何连通。这来自一个更精确的事实：若 $E$ 自由且秩为 $r\geq n$，则 $\operatorname{Grass}_n(E)$ 可以由 $\binom rn$ 个开子集覆盖，每个开子集都在 $S$ 上同构于相对维数为 $n(r-n)$ 的仿射空间。\cmt{D19-11-05}{法文第61页在“局部自由”后划去了“有限型”，但紧随的坐标开覆盖及随后提升截面的论述都在有限秩情形使用。正文明确保留此处已支持的有限秩范围，不将其推广至任意无限秩自由模。原稿亦未说明 $n>r$ 的边界；此时代表预概形为空。非空几何连通纤维断言因此限于 $n\leq r$，空部分仍在 $S$ 上光滑。}
这一覆盖来自于借助 $E$ 的一组基选取 $\binom rn$ 个坐标直和分解，即分裂正合列。\cmt{D19-11-06}{法文第61页的首项确为 $E'$，并在紧接的下一行说明它的秩为 $n$；英文第41页把正合列的首项误排成 $E$。此处沿用旧中文的正确公式，并将错误来源限定为英译。}
\[
(s)\qquad0\longrightarrow E'\longrightarrow E
\longrightarrow E''\longrightarrow0,
\qquad\operatorname{rank}E'=n.
\]

这样的正合列确定 $\operatorname{Grass}_n(E)$ 的一个子函子 $\operatorname{Grass}_n(s)$：只取那些使复合同态
\[
E'_{S'}\longrightarrow E_{S'}\longrightarrow G
\]
为满射、从而为同构的局部自由秩 $n$ 商 $G$。包含态射
\[
\operatorname{Grass}_n(s)\longrightarrow\operatorname{Grass}_n(E)
\]
可由开浸入表示；另一方面，在上述选定的坐标分裂下，$\operatorname{Grass}_n(s)$ 典范同构于向量丛。\cmt{D19-11-07}{法文第61页没有英译第42页重复出现的“可由开浸入表示；另一方面”，并且子函子的下标使用正合列名 $s$；英文混写了 $s$ 与 $S$。这里的开图参数化从 $E$ 到 $E'$ 的、在 $E'$ 上为恒等的同态；选定分裂后，它由 $E''\to E'$ 给出。在本书 $\mathbf V(M)=\operatorname{Spec}\operatorname{Sym}M$ 的约定下，下式正好代表这种同态。若只给出任意未选分裂的短正合列，则开图是相应向量丛的挠子，并无无条件的典范零截面。旧注D17-E06的这一数学限制继续保留。}
\[
\mathbf V\bigl(\mathscr{H}\!om_{\mathcal O_S}(E',E'')\bigr).
\]

这种特殊结构的一个推论是：若 $s\in S$，则在 $E$ 局部自由、秩有限的条件下，$\operatorname{Grass}_n(E)$ 的每个 $k(s)$-值点都可提升为 $s$ 的某个邻域上的截面。另一方面，若 $S=\Spec k$，$k$ 是无限域，则 $\operatorname{Grass}_n(E)$ 的每个非空开子集都包含一个 $k$-有理点。

$\operatorname{Grass}_n(E)$ 的一个 $S$-值点，即 $E$ 的一个局部自由秩 $n$ 商模 $G$，典范地确定 $\mathbf P(E)$ 的一个闭子预概形 $\mathbf P(G)$。这样的子预概形称为 $\mathbf P(E)$ 的\emph{线性子簇}（可能混淆时，说相对于 $S$ 的线性子簇）；这里在定义名称时不限定或指定 $G$ 的秩。\cmt{D19-11-08}{法文第62页已经在“这样的子预概形”上方手写补入“但不规定 $G$ 的秩”。英译将此括注标成译者说明，旧中文因此写作“英译者说明”；本次恢复其法文手稿已有的身份。为容纳下面的零商约定，记 $\operatorname{Grass}_0(E)=S$。}
若 $n\geq1$，它是相对维数 $n-1$ 的射影纤维丛；若 $n=0$，它为空。立即可以验证，若知道对应的线性子簇，就知道 $\operatorname{Grass}_n(E)$ 的截面，也就是商 $G$。这样，对于 $n\geq1$，Grassmann 概形可解释为函子“$\mathbf P(E)_{S'}$ 的相对维数 $n-1$ 的线性子簇”的代表，其中 $S'$ 变动。

此外，还能给出后一函子、也就是相对维数 $m$ 的线性子簇概念的内在刻画：它们是 $P$ 的闭子预概形，在 $S$ 上光滑，处处具有相对维数 $m$，并且在每个 $s\in S$ 处“射影次数为一”。这一刻画将在后面的章节中给出，我们这里完全用不到它。\cmt{D19-11-09}{法文第62页明确写有“$P$ 的闭子预概形，在 $S$ 上光滑，处处相对维数 $m$”这一整段；英文第42页漏去了这些条件，旧中文又据英译把它判作“底本语法有缺口”。本次按法文补回，撤销对法文缺失这些条件的判断。此段仍是作者的预告，本节没有给出这项刻画的证明；不能把本次补回原文称作已独立完成其一般性证明。}

再假设 $E$ 局部自由、秩为 $r$，并以 $E^\vee$ 表示其对偶。由极性对应，得到典范同构
\[
\operatorname{Grass}_n(E)\simeq
\operatorname{Grass}_{r-n}(E^\vee),
\]
它把 $E$ 的商 $G$ 对应到 $E^\vee$ 的商 $E^\vee/G^\vee$。从线性子簇的观点看，$P$ 中相对维数为 $m$ 的线性子簇 $L$，对应到 $P^\vee$ 中相对维数为
\[
r-m-2=(r-1)-1-m
\]
的对偶线性子簇 $L^0$，也就是说，其在 $P^\vee$ 中的相对余维为 $m+1$。〔附记：这里 $r-1$ 是 $P$ 和 $P^\vee$ 在 $S$ 上的共同相对维数。〕\cmt{D19-11-10}{法文第62页打字的维数为 $(r-1)-1-m$，并手写补明相对余维 $m+1$；英文第42页多出相邻的“$+$、$-$”符号。正文依照法文整理。这里 $0\leq n\leq r$；当 $m=r-1$ 时对偶线性子簇为空，依本段约定赋予维数 $-1$。}
几何上可以如下理解这个对应。当 $r\geq1$ 时，先取 $n=r-1$，便得到同构
\[
\mathbf P(E^\vee)\simeq\operatorname{Grass}_{r-1}(E),
\]
从而可将 $P^\vee$ 的 $S$-值点（例如取 $S$-值点）认同为 $P$ 中余维为一的线性子簇，仍称为 $P$ 的\emph{超平面}。

这说明 $L^0$ 由包含 $P$ 中线性子簇 $L$ 的超平面组成；精确地说，$L^0$ 的 $S'$-值点就是 $P_{S'}$ 中包含 $L_{S'}$ 的超平面。\cmt{D19-11-11}{法文第63页括号内确有上标，写的是 $L^0$ 的 $S'$-值点；英文第42页漏成 $L$。旧中文补上上标是正确的，此处将错误来源进一步限定为英译。}
这来自下面的事实；它本应与“$P$ 的线性子簇 $L$ 确定其所来自的局部自由商 $G$”这一事实同时出现\ednote{英文此处附“$G$ 或 $Q$，字迹难辨”。这里依同段的商模记号写作 $G$。}：若 $G$ 与 $G'$ 是 $E$ 的两个局部自由商（不一定同秩），而 $G'$ 是 $G$ 的商，则
\[
\mathbf P(G')\subset\mathbf P(G)
\]
是 $\mathbf P(E)$ 的线性子簇之间的包含；这个包含正是由 $G\twoheadrightarrow G'$ 诱导的态射。\cmt{D19-11-12}{法文第63页右端连接词受到扫描边缘影响；英译第43页明确采用“if”。正文保留这一充分条件形式。旧中文把它写成“等价于”；在本段局部自由有限秩的范围内，逆向也能由相容的普遍商及 $p_*\mathcal O_{\mathbf P(G')}(1)\simeq G'$ 证明，但这是单列的译校说明，不冒作此次已经辨出的法文连接词。}

以上是我们至少应当掌握的连锁事实〔法文：\textit{sorites}〕\ednote{请询问 A.G. 或 Deligne，最好用哪个词。〔英文底本脚注59。〕}。在本段\ednote{A.G. 所说的 paragraph 是什么意思？〔英文底本脚注60。〕}所有其他编号小节写出之前，无论如何还不能确定完整的清单。

我觉得也宜引入函子 $\operatorname{Grass}(E)$，其在 $S'$ 上的值为 $E_{S'}$ 的局部自由、秩不预先指定的商模所成的集合。\ednote{英文在这里附“字迹不清，请询问 A.G.”；法文第63页手写的“局部自由、秩不预先指定”可辨，正文据此译出。}
于是它由下述不交并代表。\transnote{这里“秩不预先指定”指局部有限秩，秩函数可在不同开闭部分取不同的非负整数；不是允许无限秩商。否则下面的有限整数指标不交并不能代表那个更大的函子。法文这一手写补记位于一个已划去的旧§12标题之后，而有效§12标题另起于法文第64页；据其内容及英文第43页的编排，将补记留在§11末。被划去的旧标题可辨为线性截面有关结果的“补叙”，中间覆打词仍未可靠辨定。}
\[
\coprod_{n\geq0}\operatorname{Grass}_n(E)
\]
$\mathbf P(E)$ 的线性子簇确实由 $\operatorname{Grass}(E)$ 在 $S$ 上的截面确定。〔附记：若 $S$ 不连通，秩、也就是相对维数，可以变化。〕\ednote{英文在末句附“此处字迹略有不清，请与 A.G. 核实”。法文第63页可辨的末句到“若 $S$ 不连通”为止，手写左方括号没有可见的对应右括号；正文仅闭合排版括号。}
